\section{Revisão da literatura}
\label{sec:literatura}

\subsection{A Chapada do Apodi: o dano documentado}

O conhecimento sobre a Chapada vem da saúde coletiva. O estudo coordenado por
Raquel Rigotto acompanhou a região entre 2007 e 2010 e reuniu, em livro,
medições de resíduos na água, inquérito com trabalhadores e a observação direta
dos voos \citep{rigotto2011livro}. Dele saem os dados de operação usados na
Seção~\ref{sec:politica} \citep{marinho2011dimensao, teixeira2011rebulico}. Um
inventário independente dos ingredientes ativos comercializados no Baixo
Jaguaribe classifica a maioria como de alto ou médio potencial de contaminação
de águas superficiais \citep{gama2013inventario}.

Dois trabalhos tratam de desfechos de saúde. \citet{rigotto2013trends} comparam
Limoeiro do Norte, Quixeré e Russas com doze municípios de agricultura familiar
entre 2000 e 2010: a taxa de internação por neoplasia é 1,76 vez maior e a de
óbito por neoplasia, 1,38 vez; não há diferença em óbito fetal nem em
malformação. \citet{aguiar2017} estuda oito famílias da comunidade do Tomé com
casos de malformação congênita e puberdade precoce, sete delas com nascimento
ou diagnóstico depois de 2014. São estudos descritivos, sem grupo de comparação
que sustente inferência causal, e os autores o reconhecem. Fornecem a este
projeto o problema e a comparação que o motiva, entre a fruticultura patronal e
a agricultura familiar.

\subsection{Agrotóxicos e saúde no Brasil: proxies de uso}

A referência com identificação causal é \citet{dias2023down}. Os autores
exploram a liberação da soja transgênica, a aptidão agroclimática dos
municípios e a direção dos rios, e encontram aumento da mortalidade infantil a
jusante das áreas que mais expandiram o uso de glifosato. O químico e o canal
são outros: herbicida em grãos, por água de superfície.

Dois trabalhos do programa de economia aplicada da Universidade Federal de
Viçosa tentam medir efeitos em escala nacional com dados municipais.
\citet{rodrigues2022tese} mede a intensidade de uso pela despesa com
agrotóxicos por hectare, nos Censos Agropecuários de 2006 e 2017, e a relaciona
às intoxicações notificadas, com a temperatura como instrumento; num segundo
ensaio, estima o valor do produto marginal do agrotóxico e conclui por uso
excessivo em quase todo o país, com a exceção do Nordeste. \citet{januzzi2025}
constrói um índice de exposição a partir da intensidade de uso por cultura em
Mato Grosso e da área plantada de sete culturas em cada município, e não
encontra efeito sobre baixo peso, prematuridade ou mortalidade infantil nas
especificações com efeitos fixos de município e de ano. Um terceiro trabalho do
grupo estima uma função dose-resposta entre intensidade de uso e intoxicações
por escore de propensão generalizado \citep{rodrigues2022intoxicacoes}.

Três limites dessa literatura definem o espaço deste projeto. As proxies medem
quantidade ou gasto, e não método de aplicação. O índice por cultura não inclui
banana nem melão, e por isso não enxerga a fruticultura irrigada do Nordeste. E
nenhum dos trabalhos avalia uma política: a variação vem de diferenças de uso
entre municípios, com as dificuldades de identificação que isso traz. Nas
bibliografias dos três não há um estudo que aplique diferenças em diferenças a
agrotóxicos.

O mesmo conjunto de trabalhos traz dois avisos metodológicos. O primeiro é que
a notificação de intoxicação depende da infraestrutura de vigilância: o número
de municípios que notificam quase triplica entre 2006 e 2017, e a maior parte
dos óbitos por agrotóxico é de suicídios \citep{rodrigues2022tese}. O segundo é
que, na janela de 2019 a 2022, a presença ou ausência do efeito fixo de ano
decide a significância do resultado de \citet{januzzi2025}. Nesse período, só um
grupo de comparação contemporâneo separa efeito de choque comum.

\subsection{Evidência com identificação causal em outros países}

\citet{calzada2023bananas} estudam a fumigação aérea de bananais no Equador
com o endereço da mãe, o perímetro das plantações e o registro dos voos. A
estratégia é de diferenças em diferenças sazonal: comparam mães que moram perto
e longe dos bananais, em gestações que atravessam ou não a estação de fumigação
intensa. Entre as que moram a até 150 metros, a exposição intensa reduz o peso
ao nascer em 80 gramas (erro-padrão de 31) e eleva a chance de baixo peso; o
efeito está no primeiro e no segundo trimestre de gestação e desaparece além de
150 metros. É o análogo deste projeto em cultura, método, classe química e
desfecho, e a fonte do efeito de referência do cálculo de poder.

\citet{camacho2017health} usam a aspersão aérea de herbicida sobre cultivos
ilícitos na Colômbia, medida por município e dia, e encontram aumento de
consultas dermatológicas e respiratórias e de perdas gestacionais. Interessa
aqui por ter exposição municipal, como a deste projeto, e por mostrar a perda
gestacional como desfecho sensível. \citet{larsen2017agricultural} encontram
efeitos sobre nascimentos na Califórnia apenas na cauda superior da
distribuição de exposição, o que antecipa a dificuldade de detectar efeitos em
médias.

\subsection{Custo social e instrumentos}

\citet{soares2009socialcost, soares2012impactos} estimam o custo da intoxicação
aguda no Paraná e o comparam ao benefício produtivo do agrotóxico. Fornecem o
método de valoração por custo da doença. \citet{cunha2020incentivos} calculam
em R\$~9,86 bilhões a desoneração fiscal dos agrotóxicos em 2017: no Brasil, o
instrumento de preço opera como subsídio. A tradição que organiza a escolha
entre proibir e taxar vem de \citet{weitzman1974} e, para agrotóxicos, de
\citet{lichtenberg1986damage} e \citet{lichtenberg1988}.

\subsection{Método}

O estimador é o de \citet{callaway2021multiple}, com a versão duplamente
robusta de \citet{santanna2020drdid}. A leitura dos períodos anteriores ao
tratamento segue \citet{roth2023trending} e \citet{rambachan2023credible}. A
diferença tripla segue \citet{olden2022triple}. Para inferência com poucas
unidades tratadas, \citet{conley2011} e \citet{ferman2019few}. O vocabulário de
avaliação de políticas (teoria da mudança, validade, poder) é o de
\citet{gertler2016}.
